\chapter{Concurrency I --- \texttt{jthread} and Cooperative Cancellation}
\label{ch:c20-jthread}

\begin{quote}
\emph{\texttt{std::jthread} fixes the two most dangerous defaults of \texttt{std::thread}: it joins automatically in its destructor (no more \texttt{std::terminate} from a forgotten \texttt{join}), and it carries built-in cooperative cancellation via \texttt{std::stop\_token}. With \texttt{std::stop\_source} and \texttt{std::stop\_callback}, C++20 gives the standard library a structured, RAII-clean way to start, stop, and wait for threads.}
\end{quote}

\section{The std::thread Problem jthread Solves}
\label{sec:c20-jthread-problem}

\begin{lstlisting}[language=C++,caption={The std::thread footgun vs jthread}]
#include <thread>

void with_plain_thread() {
    std::thread t{[]{ /* work */ }};
    // Returning without join()/detach() calls std::terminate().
    t.join();   // mandatory, easy to forget
}
void with_jthread() {
    std::jthread t{[]{ /* work */ }};
    // ~jthread() requests stop and joins automatically.
}
\end{lstlisting}

\section{Automatic Joining: RAII for Threads}
\label{sec:c20-jthread-raii}

\begin{lstlisting}[language=C++,caption={Scope-bounded thread lifetime}]
#include <thread>
#include <vector>

void process(const std::vector<int>& data) {
    std::jthread worker{[&]{ for (int x : data) { /* ... */ } }};
    if (data.empty()) return;   // ~jthread joins here — no leak, no terminate
}                               // ~jthread joins here on the normal path too
\end{lstlisting}

\section{The Cancellation Trio}
\label{sec:c20-jthread-trio}

\begin{center}
\begin{tabular}{ll}
\hline
\textbf{Type} & \textbf{Role} \\
\hline
\texttt{std::stop\_source} & requester: \texttt{request\_stop()} sets the state \\
\texttt{std::stop\_token} & observer: polls \texttt{stop\_requested()} \\
\texttt{std::stop\_callback} & reaction: runs a callable on stop \\
\hline
\end{tabular}
\end{center}

\begin{lstlisting}[language=C++,caption={How the three types share one stop-state}]
#include <stop_token>

std::stop_source source;
std::stop_token  token = source.get_token();
bool before = token.stop_requested();   // false
source.request_stop();
bool after  = token.stop_requested();   // true
\end{lstlisting}

\section{Writing a Cancellable Thread Function}
\label{sec:c20-jthread-cancellable}

\begin{lstlisting}[language=C++,caption={A cooperatively cancellable worker}]
#include <thread>
#include <stop_token>
#include <chrono>

void worker(std::stop_token st, int id) {   // stop_token MUST be the first param
    while (!st.stop_requested()) {
        std::this_thread::sleep_for(std::chrono::milliseconds(100));
    }
}
int main() {
    std::jthread t{worker, 42};   // jthread injects the stop_token as arg 0
    std::this_thread::sleep_for(std::chrono::seconds(1));
}   // ~jthread() calls request_stop() then join()
\end{lstlisting}

There is no forced thread termination in C++; the thread must voluntarily poll at safe points.

\section{stop\_source: Requesting a Stop from Outside}
\label{sec:c20-jthread-source}

\begin{lstlisting}[language=C++,caption={Explicit and external stop requests}]
#include <thread>
#include <stop_token>

std::jthread t{[](std::stop_token st){ while (!st.stop_requested()){} }};
t.request_stop();                 // non-blocking; thread stops at next poll

std::stop_source group;           // fan-out cancellation
auto tok = group.get_token();
std::jthread b{[tok]{ while (!tok.stop_requested()){} }};
group.request_stop();             // stops every worker holding a token from 'group'
\end{lstlisting}

\section{stop\_callback: Reacting to Cancellation}
\label{sec:c20-jthread-callback}

\begin{lstlisting}[language=C++,caption={Running code the instant a stop is requested}]
#include <stop_token>
#include <atomic>

void register_reaction(std::stop_token st, std::atomic_flag& wake) {
    std::stop_callback cb{st, [&wake]{
        wake.test_and_set();
        wake.notify_one();
    }};
    // Fires immediately if already stopped, else on request_stop().
}
\end{lstlisting}

The callback runs on the thread calling \texttt{request\_stop()}; \texttt{cb}'s destructor deregisters it and blocks until any in-progress call finishes.

\section{Interrupting a Blocking Wait}
\label{sec:c20-jthread-wait}

\begin{lstlisting}[language=C++,caption={A condition-variable wait that honors cancellation}]
#include <stop_token>
#include <condition_variable>
#include <mutex>
#include <queue>

template <typename T>
class BlockingQueue {
    std::mutex m_;
    std::condition_variable_any cv_;
    std::queue<T> q_;
public:
    bool pop(std::stop_token st, T& out) {
        std::unique_lock lock{m_};
        if (!cv_.wait(lock, st, [&]{ return !q_.empty(); }))
            return false;             // stop requested -> bail out
        out = std::move(q_.front());
        q_.pop();
        return true;
    }
};
\end{lstlisting}

Requires \texttt{condition\_variable\_any}; the overload registers an internal \texttt{stop\_callback} that wakes the waiter on \texttt{request\_stop()}.

\section{Professional Insights}
\label{sec:c20-jthread-insights}

\textbf{Default to \texttt{std::jthread} for every new thread} --- automatic join eliminates the forgotten-\texttt{join} \texttt{terminate} crash; reserve \texttt{std::thread} for interop.

\textbf{Cancellation is cooperative} --- design frequent, responsive checkpoints; an unchecked loop hangs the destructor's join.

\textbf{Make blocking waits stop-aware} with the \texttt{condition\_variable\_any} + \texttt{stop\_token} overload, or the worker sleeps through cancellation.

\textbf{Use one shared \texttt{stop\_source} for fan-out cancellation} of a pool --- a single \texttt{request\_stop()} shuts down the group race-free.

\textbf{Rely on \texttt{stop\_callback}'s destructor semantics} for safe cleanup, but keep callbacks short and non-blocking (signal, don't compute).
